\section{Shrinking spheres and harmonic transfer}
\label{sec:spheres}\label{sph:section}

We now work with the contact data of Proposition~\ref{bdy:contact}.
The two systems agree on the lower side of a smooth hypersurface
through the origin, and their Green matrices agree when both arguments
are on that side.  Our aim is to represent a map between their local
harmonic sections by integration over small spheres.  Product
continuation constructs this map at each positive radius.  A separate
estimate in the next section will control its density as the radii
collapse.

The sphere geometry and radius profiles below are adapted from
\cite[Section~4]{ScalarCompanion2026}.  We give their proofs, including
the strict Hessian calculation, and formulate the resulting transfer
for the present matrix equations.  The coordinates need not be harmonic.

\subsection{A concave depth and the radius profiles}

Apply the spatial dilation~\eqref{bdy:dilation-formulas}.  In this section we
suppress the dilation subscript on $g_j,A_j,G_j,F_j$ and on the paired
functions, while retaining $\lambda$ for the dilation parameter.
On every fixed bounded coordinate region, $g_j\to I$ smoothly and
$A_j\to0$ smoothly as $\lambda\downarrow0$.  The harmonic frames and
their inverses have uniform smooth bounds there.  All bounded regions
and positive separation scales used below are fixed before the final
upper bound on $\lambda$ is chosen.

Put
\[
 Y=\{y=(y',y_n):|y'|\le1,\ -1\le y_n\le1\},
 \qquad t_0(y)=y_n+|y'|^2.
\]
For a sufficiently large fixed $L$, define
\begin{equation}\label{sph:normalization}
 \gamma=e^{-2Lt_0}g_1,\qquad a=\gamma^{-1},\qquad
 Q=\frac{n g_2^{-1}}{\operatorname{tr}(\gamma g_2^{-1})},
 \qquad h=Q-a.
\end{equation}
These tensors are defined on a fixed neighborhood of $Y$, also used
for the first-variable points.  Specify the scalar multipliers in the
harmonic-frame convention~\eqref{bdy:harmonic-normalization} by
\begin{equation}\label{sph:frame-normalization}
 \mathcal L_jv=-c_jF_j^{-1}L_j(F_jv),\qquad
 c_1=e^{2Lt_0},\qquad
 c_2=\frac{n}{\operatorname{tr}(\gamma g_2^{-1})}>0.
\end{equation}
Their second-derivative coefficient matrices are $a$ and $Q$,
respectively, and their zeroth-order terms vanish because the columns
of $F_j$ are harmonic.  These changes leave the solution spaces
unchanged after the frame identification.  For product continuation
we use the positive-principal-symbol convention: $-\mathcal L_j$ in
harmonic frames, or equivalently $c_jL_j$ in unitary frames.  The
corresponding principal metrics are $\gamma$ and $Q^{-1}$.
The normalization ensures
\begin{equation}\label{sph:trace-free}
 \operatorname{tr}(\gamma h)=0,\qquad h\longrightarrow0
 \quad\hbox{in every fixed smooth norm as }\lambda\downarrow0.
\end{equation}

Define a depth function by
\begin{equation}\label{sph:depth}
 t(y)=\int_{1/8}^{t_0(y)}e^{-2Ls}\,ds.
\end{equation}
It is negative near the origin and has nonvanishing differential.
Its useful feature is strict concavity for the reference metric:
\begin{equation}\label{sph:depth-hessian}
 \operatorname{Hess}_{\gamma}t
 =e^{-2Lt_0}\bigl(
     \operatorname{Hess}_{g_1}t_0
       -L|dt_0|_{g_1}^2g_1\bigr)\le-c\gamma.
\end{equation}
Indeed the conformal connection formula cancels the two
$dt_0\otimes dt_0$ terms.  The norm of $dt_0$ is bounded below on the
fixed region, whereas $\operatorname{Hess}_{g_1}t_0$ is uniformly bounded
for small $\lambda$.  Choosing $L$ first proves the last inequality.
All constants describing this reference geometry are henceforth fixed.

For an integer $p_*\ge2$, an amplitude $R_*>0$, and $0<\delta\le1$, put
\begin{equation}\label{sph:radius}
 \ell_\delta(t)=\delta\log(1+e^{t/\delta}),\qquad
 r_\delta(y)=R_*\ell_\delta(t(y))^{p_*},\qquad
 f_*=\log r_\delta,\qquad \beta=|df_*|_\gamma.
\end{equation}
We use the $\gamma$-geodesic ball with center $y$ and radius
$r_\delta(y)$.  These balls shrink to a point on $\{t\le0\}$, including
a neighborhood of the contact point.  At positive depths bounded away
from zero, their radii retain a positive lower bound.  This will let us
place cutoff derivatives where the transfer is already controlled.

\begin{lemma}[Uniform profile bounds]\label{sph:profiles}
For the geometry just fixed and all sufficiently small $\lambda$, the
profiles~\eqref{sph:radius} satisfy
\begin{equation}\label{sph:profile-bounds}
 \beta\ge cp_*,\qquad
 \operatorname{Hess}_\gamma f_*\le-c\beta\gamma,
 \qquad |\partial^j f_*|\le C_j\beta^j\quad(j\ge1),
\end{equation}
uniformly in $0<\delta\le1$.  The constants $c,C_j$ are independent of
$p_*\ge2$.  For fixed $p_*,R_*$, the functions $r_\delta^{1/p_*}$ have
a uniform Lipschitz bound.  The radii increase strictly with $\delta$,
converge to zero where $t\le0$, and obey
\begin{equation}\label{sph:amplitude-smallness}
 \sup_{0<\delta\le1,\ y\in Y}r_\delta\beta^2
 \longrightarrow0\quad\hbox{as }R_*\downarrow0
\end{equation}
for each fixed $p_*$.
\end{lemma}

\begin{proof}
Write $a_\delta=(\log\ell_\delta)'$.  With $s=t/\delta$,
\[
 a_\delta(t)=\delta^{-1}A(s),\qquad
 A(s)=\frac{e^s}{(1+e^s)\log(1+e^s)}.
\]
The function $A$ is positive and decreasing: setting $u=e^s$ gives
\[
 A'(s)=\frac{u\bigl(\log(1+u)-u\bigr)}
              {(1+u)^2\log(1+u)^2}<0.
\]
At the negative end $A$ tends to one, and at the positive end it is
comparable to $(1+s)^{-1}$.  Differentiating its convergent expressions
at the two ends gives $|A^{(j)}|\le C_jA^{j+1}$; positivity gives the
same bounds on the remaining compact interval.  Consequently, on the
fixed range of $t$,
\[
 a_\delta\ge c>0,\qquad a_\delta'\le0,\qquad
 |\partial_t^j\log\ell_\delta|\le C_j a_\delta^j,
 \qquad a_\delta\le\ell_\delta^{-1}.
\]
Since $|dt|_\gamma$ is bounded above and below,
$\beta=p_*a_\delta|dt|_\gamma$.  The identity
\[
 \operatorname{Hess}_\gamma f_*
   =p_*a_\delta\operatorname{Hess}_\gamma t
       +p_*a_\delta'\,dt\otimes dt
\]
and~\eqref{sph:depth-hessian} prove the first two bounds
in~\eqref{sph:profile-bounds}.  The chain rule gives the derivative
bounds; their constants can be independent of $p_*$ because
$p_*a_\delta^j\le p_*^j a_\delta^j$ for $j\ge1$.

The estimate $0<\ell_\delta'<1$ proves the Lipschitz assertion.
Moreover,
\[
 \partial_\delta\ell_\delta
   =\log(1+e^s)-\frac{s e^s}{1+e^s}>0.
\]
The displayed expression tends to zero at both ends and its derivative
in $s$ is $-s e^s/(1+e^s)^2$, so it is strictly positive at finite $s$.
The limit of $\ell_\delta$ is $\max(t,0)$.  Finally,
\[
 r_\delta\beta^2\le C R_*p_*^2\ell_\delta^{p_*-2},
\]
and $\ell_\delta$ is uniformly bounded on the fixed depth interval.
This proves~\eqref{sph:amplitude-smallness}.
\end{proof}

\subsection{The core, cutoff, and order of choices}\label{sph:core-cutoff}

Choose $t_b>0$ so small that $3t_b<t(1/4)$, where $t(1/4)$ denotes
the value of the increasing function in~\eqref{sph:depth} at
$t_0=1/4$.  Define
\begin{equation}\label{sph:cutoff}
 \begin{gathered}
 K=\{y\in Y:y_n\ge-1/2,\ t(y)\le2t_b\},\\
 \chi\in C_c^\infty(Y^\circ),\qquad \chi=1\text{ near }K,\\
 \operatorname{supp}(d\chi)\subset
       \{y_n<-1/4\}\cup\{t>t_b\}.
 \end{gathered}
\end{equation}
Such a cutoff exists.  Indeed, on $K$ one has $t_0<1/4$, and thus
$|y'|^2<3/4$ and $y_n<1/4$; hence $K\Subset Y^\circ$.
Multiply a cutoff changing from zero to one between $y_n=-3/4$
and $y_n=-1/2$ by a cutoff of $t$ changing from one to zero between
$2t_b$ and $3t_b$, moving its endpoints slightly to leave room around
$K$.  The inequality $3t_b<t(1/4)$ keeps the support away from the
lateral and upper faces of $Y$.  On the second cutoff region,
\begin{equation}\label{sph:positive-depth-radius}
 r_\delta\ge R_*t_b^{p_*}>0.
\end{equation}
Every point of $K$ can be joined to the known lower region by the
vertical segment with the same $y'$ and initial height $-1/2$.
That segment stays in $K$; both $t$ and $r_\delta$ are nondecreasing
along it.  For sufficiently small $\lambda$, its initial point and a
neighborhood of it lie in the known side.

We will use the following temporary metric comparison:
\begin{equation}\label{sph:smallness}
 H_8(y):=\sum_{|\alpha|\le8}|\partial^\alpha h(y)|
       \le\varepsilon r_\delta(y)\qquad(y\in Y).
\end{equation}
Coordinate norms in this definition are uniformly equivalent to the
reference-metric norms.  This condition will later be improved and
then removed by a continuity argument.

\begin{remark}[Parameter choices]\label{sph:choices}
The reference geometry, $t_b$, $K$, and $\chi$ are fixed first.
Choose $p_*$ sufficiently large and then $\varepsilon>0$ sufficiently
small.  The lower bounds on $p_*$ and upper bounds on $\varepsilon$
used below depend only on the fixed geometry and coefficient bounds,
before $R_*$ and the final dilation are chosen.
Choose $R_*$ so that every ball lies in a normal neighborhood with
fixed coordinate room and
\begin{equation}\label{sph:radius-smallness}
 r_\eta|d\log r_\eta|_\gamma^2\le\varepsilon
 \qquad(0<\eta\le1,\ y\in Y).
\end{equation}
The fixed-separation scales determined by these choices are positive,
although they may be small.  Only then reduce $\lambda$ to meet the
fixed-separation initialization and the coefficient bounds.  The
parameter $\delta$ remains free subject to~\eqref{sph:smallness}.
In particular the upper bound on $\lambda$ does not depend on
$\delta$.
\end{remark}

\subsection{The strict Hessian test for a moving sphere}

The difficulty in approaching the diagonal is that the radius varies
with the center.  The leading terms from its gradient cancel in the
product Hessian test.  The strict concavity of its logarithm provides
the positive remainder.

\begin{lemma}[Strict sphere geometry]\label{sph:strict}
Let $r=r_\delta$ satisfy~\eqref{sph:profile-bounds},
\eqref{sph:smallness}, and~\eqref{sph:radius-smallness}.
For $p_*$ sufficiently large and $\varepsilon$ sufficiently small,
the hypersurface
\[
 \rho(x,y)=\tfrac12\bigl(d_\gamma(x,y)^2-r(y)^2\bigr)=0
\]
satisfies the strict criterion of Proposition~\ref{cross:criterion}
for principal metrics $\gamma$ in the first factor and $Q^{-1}$ in
the second.  Both partial gradients are nonzero, and the positive
side is the exterior of the moving ball.
\end{lemma}

\begin{proof}
Put $b=\nabla^\gamma r$, so that $|b|=r\beta$.  Parallel transport
along the radial $\gamma$-geodesic identifies the endpoint tangent
spaces; let $\omega$ be its unit direction from $y$ to $x$.
Multiply the normalized vectors in Proposition~\ref{cross:criterion}
by the common factor $r$.  Their real and imaginary parts can be
written as
\[
 U_1=\omega+i v,\qquad v\perp\omega,\quad |v|=1,
 \qquad U_2=q+i w.
\]
All lengths in the following calculation use $\gamma$.  Set
$s=\omega+b$.  In a $\gamma$-orthonormal frame at $y$, the exact
conditions on the second vector are
\begin{equation}\label{sph:test-vectors}
 q=\frac{Qs}{s^TQs},\qquad s\cdot w=0,\qquad
 w^TQ^{-1}w=\frac1{s^TQs}.
\end{equation}
Here the matrix of $a$ is the identity.  Therefore
\begin{equation}\label{sph:test-vector-bounds}
 s\cdot U_2=1,\qquad
 q=\frac{s}{|s|^2}+O(|h|),\qquad
 |w|=|s|^{-1}+O(|h|).
\end{equation}
The assumptions make $b$ and $h$ small.  Both partial gradients of
$\rho$ are consequently nonzero, and $|U_2|$ is bounded above and below.

With the product $\gamma$-connection, the endpoint Hessian of half
the squared distance has value
\begin{equation}\label{sph:distance-hessian}
 |U_1-U_2|^2+O(r^2)(|U_1|^2+|U_2|^2).
\end{equation}
To obtain the error, parametrize the radial geodesic on $[0,1]$.
In a parallel frame its Jacobi field with prescribed endpoint values
differs from the affine interpolant by $O(r^2)$ in $C^1$, because the
curvature is bounded and the velocity has size $r$.  The second
variation formula gives~\eqref{sph:distance-hessian}.
Replacing the second connection by that of $Q^{-1}$ contributes only
$O(r^2)$: its difference from the $\gamma$-connection is
$O(|h|+|\partial h|)=O(\varepsilon r)$, while $|d\rho|=O(r)$.

The relation $s\cdot U_2=1$ implies
$\omega\cdot(U_1-U_2)=b\cdot U_2$.  Also
\[
 \operatorname{Hess}_\gamma(r^2/2)
   =2\,dr\otimes dr+r^2\operatorname{Hess}_\gamma f_*.
\]
Thus the Hessian sum in the criterion, multiplied by $r^2$, equals
\begin{equation}\label{sph:test-expression}
 |\pi_{\omega^\perp}(U_1-U_2)|^2-|b\cdot U_2|^2
       -r^2\operatorname{Hess}_\gamma f_*(U_2,\overline{U_2})
       +O(r^2).
\end{equation}
For a real bilinear form, evaluation on a complex vector and its
conjugate means the sum of the evaluations on its real and imaginary
parts, as in the criterion.

We record the cancellation quantitatively.  The real projected
difference is
$-\pi_{\omega^\perp}b+O(|b|^2+|h|)$.  The imaginary part obeys
\[
 |\pi_{\omega^\perp}w|
    =1-\omega\cdot b+O(|b|^2+|h|).
\]
Comparing this length with $|v|=1$ and squaring gives
\[
 |\pi_{\omega^\perp}(U_1-U_2)|^2
 \ge |b|^2-C\bigl(|b|^3+|b||h|+|h|^2\bigr).
\]
The real and imaginary directions of $U_2$ are almost orthonormal;
hence
\[
 |b\cdot U_2|^2\le
       \bigl(1+C(|b|+|h|)\bigr)|b|^2.
\]
Using $|b|=r\beta$, $|h|\le\varepsilon r$, and
\eqref{sph:profile-bounds}, expression~\eqref{sph:test-expression}
is bounded below by
\begin{equation}\label{sph:test-lower-bound}
 c r^2\beta-
 C\bigl(r^2+r^3\beta^3+\varepsilon r^2\beta+
                  \varepsilon^2r^2\bigr).
\end{equation}
Divide the error terms by $r^2\beta$.  Their ratios are bounded by
$C/\beta$, $Cr\beta^2$, $C\varepsilon$, and
$C\varepsilon^2/\beta$, respectively.  Increasing $p_*$ controls the
first, and then decreasing $\varepsilon$ controls the others by
\eqref{sph:radius-smallness}.  This proves strict positivity.
\end{proof}

\begin{proposition}[Continuation to the sphere family]
\label{sph:continuation}
Make the choices of Remark~\ref{sph:choices}.  For all sufficiently
small $\lambda$ and every $0<\delta\le1$ satisfying
\eqref{sph:smallness}, the mixed Green matrix extends to a neighborhood
of
\[
 \{(x,y):y\in Y,\ d_\gamma(x,y)=r_\delta(y)\}.
\]
It solves the first system in $x$ and the conjugate second system in
the second tensor index.  On the lower center region $y_n<-1/4$ it
agrees with $G_1$, and on the prescribed positive-depth region
$t>t_b$ it agrees with the fixed-separation continuation of
Proposition~\ref{cross:initial}.  The extension is smooth for each
positive $\delta$.
\end{proposition}

\begin{proof}
Deform $\eta$ from $1$ down to $\delta$, using $(y,\omega)$ in the
unit sphere bundle of $\gamma$ as base coordinates and
\[
 x=\exp_y^\gamma(r_\eta(y)\omega).
\]
The base over the closed cylinder $Y$ is compact.  The radius has
strictly positive derivative in $\eta$, so these variables give
cylinder coordinates on every fixed interval $\delta\le\eta\le1$.
The positive side of a level is the direction of increasing $\eta$.
Since $r_\eta\ge r_\delta$, condition~\eqref{sph:smallness} holds
at every intermediate radius.  Lemma~\ref{sph:strict} applies there.

At $\eta=1$ the separations have a positive minimum on $Y$.
The same is true throughout $t>t_b$ by
\eqref{sph:positive-depth-radius}.  Proposition~\ref{cross:initial}
therefore supplies the top and positive-depth data, with open room and
uniform bounds at any fixed order.  In the lower strip $y_n<-1/4$,
use the original kernel $G_1$; this strip lies in the dilated known
side for small $\lambda$.  The initialization agrees with this kernel
on their overlap.

These two base regions cover every lateral face of $Y$: its bottom
face is in the lower strip, and on a remaining lateral or upper face
outside that strip one has $t_0\ge3/4$.  Both regions are independent
of $\eta$ and hence upward closed.  Together with a neighborhood of
the top level they supply exactly the seed required by
Lemma~\ref{cross:propagation}.  That lemma and the strict sphere test
extend the kernel through the whole cylinder, preserving agreement
on the seed.  Normal-neighborhood room and positive radii on each
fixed interval ensure that all surfaces remain off the diagonal.
\end{proof}

The proposition gives existence at every admissible positive radius.
It does not give a bound for the kernel as $\delta\downarrow0$.
We next pass from the kernel to its action on harmonic sections;
this action has a sphere density that can be estimated separately.

\subsection{The transfer and its matrix density}

Fix $\delta>0$ satisfying~\eqref{sph:smallness}, write $r=r_\delta$,
and denote the continued kernel by $\mathcal G(x,y)$.
For a first-system harmonic section $u$ defined near the closed
$\gamma$-ball centered at $y$, set
\begin{equation}\label{sph:flux}
 S(y;u)=\int_{\partial B_\gamma(y,r(y))}
 \left[\mathcal G(x,y)^*\nabla^{A_1}_{\nu_1}u(x)
       -\bigl(\nabla^{A_1}_{\nu_1,x}\mathcal G(x,y)\bigr)^*u(x)\right]
 \,dS_1(x).
\end{equation}
Here $\nu_1$ and $dS_1$ are the outward unit normal and hypersurface
density for the actual metric $g_1$.  The derivative of the kernel
acts on its first fiber index.  Consequently the displayed integrand
is a section of the second fiber, with the adjoints in the indicated
order.  For a first harmonic-frame function $v$, define
\begin{equation}\label{sph:transfer-definition}
 T_yv=F_2(y)^{-1}S(y;F_1v).
\end{equation}
Let $S_y$ denote the unit sphere of $(T_y\mathbb R^n,\gamma_y)$ and
$d\sigma_y$ its rotation-invariant probability measure.

\begin{proposition}[Matrix harmonic transfer]\label{sph:transfer}
Under the assumptions of Proposition~\ref{sph:continuation}, there
is a smooth matrix function $\psi=\psi_\delta$ on the unit sphere
bundle over $Y^\circ$ such that
\begin{equation}\label{sph:density-representation}
 T_yv=\int_{S_y}\psi(y,\omega)
       v\bigl(\exp_y^\gamma(r(y)\omega)\bigr)\,d\sigma_y(\omega).
\end{equation}
For every $y_0\in Y^\circ$ and every first harmonic-frame function
$v$ defined near $\overline{B_\gamma(y_0,r(y_0))}$, the function
$y\mapsto T_yv$ is second harmonic-frame harmonic on a neighborhood
of $y_0$.  For each pair $(v_1,v_2)$ of first and second harmonic-frame
functions defined on the fixed coordinate neighborhoods and agreeing
on the known side,
\begin{equation}\label{sph:reproduction}
 \int_{S_y}\psi(y,\omega)\,d\sigma_y(\omega)=I,
 \qquad T_yv_1=v_2(y).
\end{equation}
\end{proposition}

\begin{proof}
The Dirichlet problem on each ball is invertible.  In the unitary
frame this follows from the positive connection energy: a zero-boundary
null solution is parallel and hence zero.  Invertibility is preserved
by the harmonic-frame change and by positive scalar multiplication of
the equation.  Its Dirichlet solution operator determines the boundary
covariant normal derivative from the boundary value.  Transpose this
boundary operator onto the smooth matrix coefficient in the first term
of~\eqref{sph:flux}, then multiply by the frame factors and pull the
surface density back under $\omega\mapsto\exp_y^\gamma(r(y)\omega)$.
This gives the smooth matrix density in
\eqref{sph:density-representation}.  Smooth dependence of the ball
Dirichlet problem on its center and positive radius gives smooth
dependence of $\psi$.

To prove the local harmonicity assertion, fix $y_0$ and $u=F_1v$.
The kernel is defined on an open neighborhood of the compact sphere
over $y_0$.  Choose a fixed slightly larger surrounding surface inside
that neighborhood and inside the domain of $u$.  After restricting
the center neighborhood, this surface surrounds all the moving
spheres and the intervening collars remain in the common domain.
Green's identity on each collar replaces~\eqref{sph:flux} by its
integral over the fixed surface.  Both first-variable equations are
homogeneous in the collar.  Applying the second-variable operator
under the fixed integral now proves that $S(y;u)$ is second-system
harmonic: the columns of $\mathcal G(x,y)^*$ solve that system.
Equation~\eqref{sph:transfer-definition} proves the assertion in
harmonic frames.  This argument applies to arbitrary local harmonic
inputs; they need not extend throughout the coordinate region.

If $y_n<-1/4$, the kernel is $G_1$ and the Green representation formula
gives $S(y;u)=u(y)$.  On that region $F_1=F_2$, so $T_yv=v(y)$.
For a fixed pair $(v_1,v_2)$, both $T_yv_1$ and $v_2(y)$ solve the
second harmonic-frame equation on connected $Y^\circ$ and agree in
the lower strip.  Unique continuation proves their equality.
Finally constant columns are harmonic in both harmonic frames and
agree there, so the same argument yields the first identity
in~\eqref{sph:reproduction}.
\end{proof}

\begin{lemma}[Bounds where the cutoff varies]\label{sph:cutoff-bounds}
After the choices in Remark~\ref{sph:choices}, for every fixed integer
$m$ the angular $H^m$ norms of $\psi_\delta$ and its first base
derivatives on $\operatorname{supp}(d\chi)$ are bounded independently
of admissible $\delta$ and sufficiently small $\lambda$.
\end{lemma}

\begin{proof}
On the lower cutoff region, $T_y$ is evaluation at the center in the
harmonic frame.  Normalize its equation to principal matrix
$a=\gamma^{-1}$, use a smooth $\gamma$-orthonormal frame to pull it
back under $z\mapsto\exp_y^\gamma(rz)$, and multiply by $r^2$.
With center and radius regarded as independent parameters,
the resulting Dirichlet systems extend smoothly to $r=0$, are uniformly
elliptic, and have the Euclidean Laplacian as limit; their lower-order
coefficients vanish at $r=0$.  The Dirichlet inverses are uniformly
bounded, either by their limiting inverse and perturbation or by the
positive energy before the change of frame.  Elliptic boundary
regularity, applied after differentiation in the parameters, shows
that the evaluation density at the center has uniform angular smooth
bounds and smooth center--radius dependence.  The center remains a
fixed positive distance from the unit sphere in these coordinates.
Substituting the varying radius costs only
$|dr|_\gamma=r\beta$, which is uniformly bounded by
\eqref{sph:radius-smallness} and the lower bound for $\beta$.

This evaluation density is the density constructed in
Proposition~\ref{sph:transfer}.  Indeed both represent the same
functional on traces of functions harmonic near the closed ball.
These traces are dense in smooth boundary data.  For a prescribed
smooth trace $\varphi(\omega)$, solve on the ball of radius
$r(1+\eta)$ with boundary value $\varphi$ transported along the radial
rays.  On pulling this problem back to the closed unit ball, its
solution $z_\eta$ converges in every $C^k$ norm to the solution $z_0$
at radius $r$, by smooth dependence of the Dirichlet inverse and
boundary regularity.  Its trace on the original sphere is
$z_\eta(\omega/(1+\eta))$, which converges in $C^k(S_y)$ to
$z_0(\omega)=\varphi(\omega)$.  The corresponding physical solutions
are harmonic near the original closed ball.  Thus the two smooth
density matrices agree.

On the other cutoff region $t>t_b$, the radii have the positive lower
bound~\eqref{sph:positive-depth-radius}; all their required derivatives
are uniformly bounded in $\delta$.  Proposition~\ref{cross:initial}
gives uniform kernel bounds there.  The smooth Dirichlet construction
of the density therefore gives the same bounds for $\psi_\delta$ and
its first base derivatives.  Compactness of the cutoff derivative
support completes the proof.
\end{proof}

We have constructed a matrix density that reproduces the matched
harmonic pairs and has controlled data wherever the cutoff changes.
The remaining estimate is a uniform bound on $\chi\psi_\delta$ over
the shrinking part of the sphere family.
